\chapter{Theoretical Background}
\label{ch:theoretical_background}

This chapter outlines the core technologies forming the SENTINEL architecture: Agentic AI, Retrieval-Augmented Generation (RAG), and Zero Trust principles.

\section{Agentic AI and LLM Quantization}
Unlike traditional Runbook Automation tools that follow rigid, pre-defined scripts, Agentic AI systems possess the autonomy to analyze contexts, evaluate evidence, and make flexible decisions based on live data streams.

However, running Large Language Models (LLMs) locally requires massive hardware resources. To make Agentic AI practically viable for enterprises, SENTINEL applies \textbf{Quantization} techniques. By reducing the precision of model weights from FP16 (16-bit) to the Q6\_K standard (approx. 6.6-bit) via the GGUF format, the VRAM requirement for the Gemma-9B model shrinks from $\approx 18$ GB to merely $\approx 7-8$ GB. This allows smooth inference on consumer-grade hardware with negligible cognitive degradation.

\section{Hybrid Retrieval-Augmented Generation (RAG)}
LLMs are inherently prone to "hallucinations" -- generating plausible but fabricated information. To counteract this, RAG forces the LLM to ground its reasoning on an authoritative knowledge base (e.g., MITRE ATT\&CK) before responding.
SENTINEL utilizes a \textbf{Hybrid Search} methodology that combines:
\begin{itemize}
    \item \textbf{Vector Search (Semantic Search - FAISS):} Captures semantic similarities in attack descriptions.
    \item \textbf{Keyword Search (BM25):} Precisely matches technical identifiers and IP addresses.
\end{itemize}
The results from both methods are merged using Reciprocal Rank Fusion (RRF), yielding optimal retrieval accuracy.

\section{Zero Trust Security and LLM Vulnerabilities}
The Zero Trust principle strictly prohibits transmitting sensitive network telemetry outside the organizational perimeter. Consequently, the entire SENTINEL architecture must operate in an Air-gapped (fully local) environment.

Furthermore, reading network logs with LLMs introduces a novel vulnerability: \textbf{Log-background Prompt Injection}. Attackers can deliberately inject deceptive control commands into the Payload (e.g., HTTP User-Agent). Without strict "Guardrails" to encapsulate the Payload (such as random tokens and delimiters), the LLM can be easily manipulated, effectively blinding or bypassing the automated SOC system.
